% from ega4-9-fr.tex
\oldpage[IV-3]{63}

On peut évidemment se borner à un complexe à trois termes de degrés
$-1$, $0$, $+1$~:
\[
\mathcal M\xrightarrow{u}\mathcal N\xrightarrow{v}\mathcal P
\]
et à $i=0$~; l'homomorphisme à considérer est alors l'homomorphisme canonique
$(\operatorname{Ker}v/\operatorname{Im}u)_s\to
\operatorname{Ker}(v_s)/\operatorname{Im}(u_s)$. Utilisant (8.9.1) et
(8.5.2, (i)), on voit qu'on peut se ramener au cas où $S$ (donc aussi $X$)
est \emph{noethérien}, et par suite $\mathcal M$, $\mathcal N$, $\mathcal P$
des $\mathcal O_X$-Modules \emph{cohérents}~; alors
$\operatorname{Im}(u)$ et $\operatorname{Ker}(v)$ sont aussi cohérents
$(\mathbf 0_{\mathrm I},5.3.4)$ et en outre il existe un voisinage $U$ de
$\eta$ tel que pour $s\in U$, on ait
$\operatorname{Ker}(v_s)=(\operatorname{Ker}(v))_s$ et
$\operatorname{Im}(u_s)=(\operatorname{Im}(u))_s$ (9.4.2)~; la conclusion
résulte alors de (9.4.2) appliqué à la suite exacte
\[
0\to\operatorname{Im}(u)\to\operatorname{Ker}(v)\to
\operatorname{Ker}(v)/\operatorname{Im}(u)\to0,
\]
compte tenu de ce que $\mathcal O_\eta=k(\eta)$ (puisque $S$ est intègre) et
par conséquent la suite
\[
0\to(\operatorname{Im}u)_\eta\to(\operatorname{Ker}v)_\eta\to
(\operatorname{Ker}v/\operatorname{Im}u)_\eta\to0
\]
est exacte.

\begin{proposition}[9.4.4]
\label{IV.9.4.4-fr}
Soient $f:X\to S$ un morphisme de présentation finie, $\mathcal F$,
$\mathcal G$, $\mathcal H$ trois $\mathcal O_X$-Modules quasi-cohérents de
présentation finie. Soient $u:\mathcal F\to\mathcal G$,
$v:\mathcal G\to\mathcal H$ deux homomorphismes de $\mathcal O_X$-Modules.
Alors l'ensemble $E$ des $s\in S$ tels que la suite
$\mathcal F_s\xrightarrow{u_s}\mathcal G_s
\xrightarrow{v_s}\mathcal H_s$ soit exacte est localement constructible.
\end{proposition}

Compte tenu de (9.2.3), il faut établir que la propriété $\mathbf P$ considérée
dans (9.4.2) est constructible. On a déjà remarqué dans la démonstration de
(9.4.2) que la condition (9.2.1, (i)) est satisfaite pour cette propriété, et
il reste à vérifier la condition (9.2.1, (ii)). Supposons donc $S$ intègre
noethérien, de point générique $\eta$, et prouvons que $E$ ou $S-E$ est un
voisinage de $\eta$. Si $\eta\in E$, notre assertion résulte de (9.4.2), et
l'on peut donc se borner au cas où $\eta\notin E$, autrement dit, la suite
$\mathcal F_\eta\xrightarrow{u_\eta}\mathcal G_\eta
\xrightarrow{v_\eta}\mathcal H_\eta$ n'est pas exacte. Distinguons alors
deux cas~:

$1^\circ$ Posons $w=v\circ u$, et supposons d'abord que
$w_\eta=v_\eta\circ u_\eta\neq0$. Comme $\mathcal F$, $\mathcal G$,
$\mathcal H$ sont cohérents, il en est de même de
$\mathcal N=\operatorname{Ker}(w)$ $(\mathbf 0_{\mathrm I},5.3.4)$~; il
résulte donc de (9.4.2) appliqué à la suite exacte
$0\to\mathcal N\to\mathcal F\xrightarrow{w}\mathcal H$ qu'il y a un
voisinage $U$ de $\eta$ dans $S$ tel que, pour $s\in S$,
$\operatorname{Ker}(w_s)=\mathcal N_s$~; en restreignant $S$, on peut donc
supposer cette relation vérifiée pour tout $s\in S$. Soit $j$ l'injection
canonique $\mathcal N\to\mathcal F$, et posons
$\mathcal M=\operatorname{Coker}(j)$~; l'exactitude à droite du foncteur
$\mathcal F\mapsto\mathcal F_s$ entraîne que
$\mathcal M_s=\operatorname{Coker}(j_s)$ pour tout $s\in S$. L'hypothèse
$w_\eta\neq0$ signifie que $\mathcal M_\eta\neq0$~; comme $\mathcal M$ est
cohérent $(\mathbf 0_{\mathrm I},5.3.4)$, il résulte de (1.8.6) qu'il y a un
voisinage ouvert $U$ de $\eta$ dans $S$ tel que $\mathcal M_s\neq0$ pour
$s\in U$, donc $w_s\neq0$ pour $s\in U$, et \emph{a fortiori} $S-E$ est un
voisinage de $\eta$.

$2^\circ$ Supposons que $w_\eta=0$~; en vertu de (8.5.2, (i)), appliqué
suivant la méthode générale de (8.1.2, \emph{a)}), il existe un voisinage
ouvert $U$ de $\eta$ tel que $w|f^{-1}(U)=0$~; remplaçant $S$ par $U$, on peut
déjà supposer $w=0$ dans $X$. Alors
$\mathcal F\xrightarrow{u}\mathcal G\xrightarrow{v}\mathcal H$ est un
complexe à trois termes $\mathcal L^\bullet$, auquel on peut appliquer
(9.4.3)~; par hypothèse on a
$(\mathcal H^0(\mathcal L^\bullet))_\eta=
\mathcal H^0(\mathcal L^\bullet_\eta)\neq0$, et
$\mathcal H^0(\mathcal L)$ est cohérent
$(\mathbf 0_{\mathrm I},5.3.4\text{ et }5.3.3)$, donc il résulte de (1.8.6)
qu'il y a un voisinage ouvert $U$ de $\eta$ tel que
$(\mathcal H^0(\mathcal L^\bullet))_s\neq0$ pour tout $s\in U$~; mais comme
on peut supposer que
$\mathcal H^0(\mathcal L_s^\bullet)=(\mathcal H^0(\mathcal L^\bullet))_s$
pour $s\in U$ par (9.4.3), on voit de nouveau que $S-E$ est un voisinage de
$\eta$ dans $S$.

